\documentclass[a4paper, oneside]{article}
\usepackage{graphicx}
\usepackage{amsthm}
\usepackage{amsmath}
\usepackage{amssymb}
\usepackage[a4paper,
            bindingoffset=0.2in,
            left=2cm,
            right=2cm,
            top=2cm,
            bottom=2cm,
            footskip=.25in]{geometry}
\usepackage[italian]{babel}
\usepackage{pgfplots}
\usepackage{tabularx}
\usepackage{tikz}
\usepackage{wrapfig}
\usepackage{color}
\usepackage[d]{esvect}
\definecolor{page}{rgb}{0.129,0.157,0.212}
\pagecolor{page}
\color{white}
\graphicspath{ {./images/} }
\usetikzlibrary{shapes.geometric}
\usetikzlibrary{datavisualization}
\usetikzlibrary{datavisualization.formats.functions}
\usetikzlibrary{patterns}
\pgfplotsset{width=10cm,compat=1.18}

\title{Appunti di Lab II}
\author{Tommaso Miliani}
\date{02-03-26}

\begin{document}
\newtheoremstyle{theoremEnv}
                {}          % Space above
                {}          % Space below
                {\slshape}  % Body font
                {}          % Indent amount
                {\bfseries} % Head font
                {.}         % Punctuation after head
                {\newline}  % Space after theorem head
                {}          % Theorem head spec
\theoremstyle{theoremEnv}

\newtheorem{definition}{Definizione}[section]
\newtheorem{theorem}{Teorema}[section]
\newtheorem{lemma}{Proposizione}[section]
\newtheorem{observation}{Osservazione}[section]
\newtheorem{corollary}{Corollario}[theorem]
\newtheorem{example}{Esempio}[section]
\newtheorem{remark}{Enunciato}[section]

\maketitle

\section{La terza equazione di Maxwell}
La seconda equazione di Maxwell era la divergenza di $\vv{E}$. Adesso, 
considerando una carica puntiforme $q$ in un sistema di riferimento,
in modo tale che non disturbi il campo elettrico,
si va a ricavare la terza equazione di Maxwell.  Supponendo 
che questa carica venga spostata da un punto $A$ ad un punto 
$B$ e si vuole determinare il lavoro del campo elettrico durante questo 
spostamento (su $\gamma$):
\begin{gather*}
    L_{\vv{F_E} } = \int_{\gamma} \vv{F_E} \cdot \vv{dl}   
\end{gather*}
Adesso la forza elettrica per la carica è
\begin{gather*}
    \vv{F_E} = \frac{Q}{4\pi\epsilon_0} \frac{\hat{r} }{r^{2}} 
\end{gather*}
Ossia $\hat{r}$ è il versore del raggio vettore che indica la carica $q$.
In coordinate sferiche, si può esprimere la carica come 
\begin{gather*}
    \vv{dl}  = dr\hat{r} + rd\theta \hat{\theta} + r\sin\theta d\phi \hat{\phi}   
\end{gather*} 
Dunque si ottiene 
\begin{gather*}
    \int_{\gamma} \vv{F_E}(dr\hat{r} + rd\theta \hat{\theta} + r\sin\theta d\phi \hat{\phi})    = \int_{\gamma} q\frac{Q}{4\pi\epsilon_0} \frac{1}{r^{2}}dr = \int_{r_A}^{r_B} q\frac{Q}{4\pi\epsilon_0} \frac{1}{r^{2}}dr 
\end{gather*}
Come si ci aspettava, non c'è lavoro per le direzioni $\hat{\theta}$ e $\hat{\phi}  $
in quanto il loro prodotto scalare con il campo è nullo. Quello che conta dunque, 
è solo l'incremento lungo $\hat{r}$. Questo ci dice che 
la forza elettrica è una forza conservativa e che, ad una certa distanza $r$ 
dall'origine del campo elettrico, se ci si muove lungo la sfera identificata da
questo raggio $r$, il campo non compie nessun lavoro. Il risultato dell'integrale è 
\begin{gather*}
    -q\frac{Q}{4\pi\epsilon_0} \left(\frac{1}{r_B} - \frac{1}{r_A}\right) = -\Delta U
\end{gather*}
Si definisce, con il risultato di questo integrale, il potenziale elettrico. Si può 
anche staccarci dalla definizione meccanica data fino ad ora e introdurre il vero potenziale 
elettrico come
\begin{align}
    \Delta V = \lim_{q  \to 0} \frac{\Delta U}{q} = \frac{Q}{4\pi\epsilon_0} \left(\frac{1}{r_B} - \frac{1}{r_A}\right)
\end{align}
Tornando all'integrale, 
\begin{gather*}
    \Delta V = \lim_{q \to 0} -\int_{\gamma} \frac{\vv{F_E} }{q}  \cdot \vv{dl} \equiv -\int_{\gamma} \vv{E} \cdot \vv{dl}   
\end{gather*}
Per una carica puntiforme dunque si è dimostrato che il lavoro non dipende dal percorso scelto. Si utilizza
ora il principio di sovrapposizione per determinare che accade per una certa distribuzione 
di carica. DAto che anche l'integrale è un operatore lineare, è dunque possibile ottenere 
una somma di integrali e dunque rimane dimostrato che tutti i campi elettrostatici sono 
\textbf{conservativi}. Il concetto che la forza è conservativa, si estende al campo
vettoriale. Dunque questa è la terza equazione di Maxwell per la forma integrale: 
\begin{align}
    V(r_B) - V(r_A) = -\oint_{r_A}^{r_B} \vv{E} \cdot \vv{dl}  \Longleftrightarrow \oint_{\gamma} \vv{E} \cdot \vv{dl} = 0  
\end{align}
Ed è indipendente da $\gamma$. L'importante è che si vada da $r_A$ a $r_B$. Un modo più diretto per 
esprimerlo è prendere due punti $A$ e $B$ e utilizzare due linee $\gamma$ da $A$ a $B$
ed una linea $\gamma'$ da $B$ ad $A$. Quello che si ottiene è che la circuitazione di $\vv{E}$ è nulla:
dunque il campo è conservativo. 

\subsection{Legare il campo elettrico al potenziale elettrico}
Partendo da un punto $A$ e andando verso un punto $A'$ che è posto a $dx, dy, dz$ di distanza
da $A$ lungo una curva $d\gamma$. SI ottiene che lo spostamento infinitesimo  mi permette di definire 
\begin{gather*}
    dV = V(A') - V(A) = -\vv{E}\cdot \vv{dl}  =  -(E_x dx + E_y dy + E_z dz) 
\end{gather*}
SI può inoltre esprimere $dV$ secondo le derivate parziali:
\begin{gather*}
    dV = \frac{\partial V}{\partial x} dx + \frac{\partial V}{\partial y}dy + \frac{\partial V}{\partial z}dz 
\end{gather*}
Nel caso in cui questa sia una funzione buona e si possa esprimere il potenziale. Uguagliando queste due quantità, e che valgono 
per qualsiasi $\vv{dl}$ si scelga, si ha che
\begin{gather*}
    E_x = -\frac{\partial V}{\partial x}  \qquad E_y = - \frac{\partial V}{\partial y}  \qquad E_z = - \frac{\partial V}{\partial z} \\
    \ \Longrightarrow \ \vv{E} = -\vv{\nabla}V     
\end{gather*} 
Dunque, a partire dalla circuitazione del campo elettrico, si può determinare il 
campo conosciuto il potenziale e che, noto il potenziale, il campo è univocamente definito.
Si deve ora trovare l'equazione di Maxwell in forma locale (ossia con il teorema 
di Stokes si deve trovare il rotore):
\begin{gather*}
    \oint \vv{E} \cdot \vv{dl} \Longleftrightarrow   \int_{S} \vv{\nabla} \times \vv{E} \cdot \hat{n} dS  = 0
\end{gather*}
Ossia il flusso del campo attraverso una superficie che 
si appoggia alla curva $\gamma$ è equivalente alla circuitazione 
del campo elettrico lungo quella curva. Questa è vera se $\vv{\nabla} \times \vv{E} = 0$. 
Ossia la \textbf{Terza equazione di Maxwell}, che deriva dalla conservatività 
della forza elettrostatica e dunque del campo elettrostatico. Si dice anche che il 
campo elettrostatico è \textbf{irrotazionale}. Dunque si dimostra che è 
uguale al suo gradiente del potenziale:
\begin{gather*}
    \vv{E} = -\vv{\nabla} \Longleftrightarrow    \vv{\nabla} \times \vv{E} =  0  
\end{gather*}
Supponendo di voler trovare
\begin{gather*}
    (\vv{\nabla} \times \vv{E}  )_1 = \frac{\partial E_3}{\partial 2}  + \frac{\partial E_2}{\partial 3} \\
\end{gather*}
Si vede intanto che il rotore di questo è nullo:
\begin{gather*}
    (\vv{\nabla} \times \vv{\nabla} V  )_1 = \frac{\partial (\vv{\nabla}V )_3}{\partial 2} + \frac{\partial (\vv{\nabla}V )_2}{\partial 3} = 0
\end{gather*}
SI dimostra anche per le altre componenti. \\ \noindent
SI definisce il rotore di $\vv{E}$ come il tensore di Ricci:
\begin{gather*}
    (\vv{\nabla} \times \vv{E}  )_i = \epsilon_{i, j, k} \frac{\partial E_k}{\partial x_j} 
\end{gather*} 
DOve si sottintende il simbolo della sommatoria. Questo tensore è una matrice tridimensionale $3x3x3$
e gode delle seguenti proprietà:
\begin{itemize}
    \item $\epsilon_{i, j, k} = 0 \Longleftrightarrow i = j, j = k, k = i$
    \item $\epsilon_{i, j, k } = 1 \Longleftrightarrow $ (1, 2, 3) sono ordinati in modo ciclico, 
    così come (3, 1, 2) o (2, 3, 1). Dunque se si fa una permutazione ciclica di una di queste, 
    il valore di questo tensore è 1.
    \item $\epsilon_{i, j, k} = -1 \Longleftrightarrow$ se si esegue una permutazione 
    non ciclica (tipo 2, 1 , 3).
\end{itemize}
Se si volesse calcolare il rotore per il primo termine posso fissare $i = 1$, dunque fisso $j = 2, 3$. Se
$j = 2$, $k = 3$, se $j = 3$, $k = 2$. Dunque di questa sommatoria a nove indici, in realtà solo due termini sono
non nulli, ossia quelli esplicitati al rotore

Dato che il potenziale è definito a meno di una costante, si impone che 
il potenziale all'infinito sia nullo. Dunque posso esprimere che
\begin{gather*}
    V(r) = \frac{Q}{4\pi\epsilon_0} \frac{1}{r}
\end{gather*}
Che riprende il concetto di messa a terra: infatti la Terra è talmente grande e talmente 
tanto carica che si può considerare che abbia potenziale nullo per qualsiasi esperimento 
che si compie su di essa. Se ci fosse una densità di carica in un volume $\tau$, si otterrebbe
un integrale di volume
\begin{gather*}
    V(r) = \int_{\tau} \frac{\rho(r')}{4\pi\epsilon_0} \frac{1}{\left| \vv{r} - \vv{r'}   \right| } d\tau  
\end{gather*}
Ogni contributo dell'integrale del potenziale ad una certa distanza $\vv{r}$ dalla carica in posizione $\vv{r'}$. Per un filo 
indefinito (o un piano) carico questa relazione non vale perché non si riesce a mettere
zero come potenziale all'infinito. Se il piano o il filo non sono infinti
invece vale. Dato che si è espresso il campo elettrico come
\begin{gather*}
    \vv{E}(r) = \frac{1}{4\pi\epsilon_0} \int_{\tau}^{}  \rho(\vv{r} ) \frac{(\vv{r} - \vv{r'} )}{\left| \vv{r} - \vv{r'}   \right| }d\tau  
\end{gather*}  
A partire da questa si definisce
\begin{gather*}
    \vv{\nabla}V = -\vv{E}  
\end{gather*}

\section{Equazione di Poisson}
Data la terza equazione di Maxwell, si prende la prima equazione:
\begin{gather*}
    \vv{\nabla} \cdot (-\vv{\nabla}V ) = \frac{\rho}{\epsilon_0}  
\end{gather*} 
Quando si fa la divergenza di un campo vettoriale, si scrive
\begin{gather*}
    \vv{\nabla} \cdot \vv{E} = \frac{\partial E_i}{\partial x_i} = -\frac{\partial }{\partial x_i} \left(- \frac{\partial V}{\partial x_i} \right) = -\frac{\partial ^{2}}{\partial x_i^{2} }V = - \frac{\partial }{\partial x_i} \frac{\partial }{\partial x_i} V  
\end{gather*}
In questa equazione, dato che compare l'indice $i$, allora si sta sommando per tutti gli $i = 1, 2, 3$. Dunque si arriva ad 
utilizzare il laplaciano per determinare l'\textbf{equazione di Poisson}:
\begin{align}
    \vv{\nabla} \cdot (-\vv{\nabla}V ) = \frac{\rho}{\epsilon_0} \Longleftrightarrow \nabla^{2}V = - \frac{\rho}{\epsilon_0} 
\end{align}
Dunque questa equazione riassume tutta l'elettrostatica e si può scrivere un'equazione 
per un campo scalare per calcolare il campo elettrico (vettoriale). Il prezzo da pagare
è che l'equazione scalare è di secondo grado. In linea di principio, qualunque problema 
elettrostatico si risolve mediante questa equazione (locale).
\begin{theorem}[Teorema di unicità (Da rivedere)]
    Supponendo di avere un volume $\tau$ nell'universo e, dentro questo 
    volume, si conosce la distribuzione di carica (dunque l'equazione di Poisson). 
    Il campo elettrico in questo volume è unico se si conoscono le 
    \textbf{boundary conditions}: ossia si devono conoscere le condizioni 
    sulla superficie S che è il bordo di $\tau$. La condizione di Diriclet mi dice che,
    se si conosce il potenziale su tutti i bordi della superficie, allora il potenziale è noto.
    Quella di ??? dice che, se si conosce il gradiente normale (ossia il campo elettrico 
    normale alla superficie), allora si conosce il potenziale.
\end{theorem}
\begin{proof}
    La dimostrazione avviene per assurdo: supponendo che non siano vere le due soluzioni del teorema, 
    ossia una soluzione $V_1(\vv{r} )$ e $V_2(\vv{r} )$ che soddisfano le ipotesi del problema, mentre o 
    sul bordo sono uguali punto per punto o si conosce il campo sulla superficie ma all'interno sono diversi. \\
    Si considera $U\vv{\nabla}U$, e si calcola la divergenza
    \begin{gather*}
        \vv{\nabla} \cdot (U\vv{\nabla}U) = \frac{\partial }{\partial x_i} \left(U \frac{\partial }{\partial x_i} U\right) = \left(\frac{\partial }{\partial x_i}U \right) \left(\frac{\partial }{\partial x_i}U \right) + U \frac{\partial }{\partial x_i}\frac{\partial }{\partial x_i}U  
    \end{gather*}
    Dunque ottengo il gradiente di $U$ al quadrato (ossia il modulo quadro del vettore 
    gradiente di $U$):
    \begin{gather*}
        \left| \vv{\nabla} U  \right|^{2} + U \nabla^{2}U  = \left| \vv{\nabla}U  \right|^{2} 
    \end{gather*}
    Dunque la divergenza di questo campo è pari al modulo quadro del campo stesso. Con questa divergenza
    si può ora fare l'integrale su $\tau$ di
    \begin{gather*}
        \int_{\tau} \vv{\nabla}(U \vv{\nabla}U ) d\tau = \int_{S} \left(U \vv{\nabla}U \cdot dS \right)   
    \end{gather*}
    Nella prima condizione questo integrale fa zero perché so che il ?? COn la seconda condizione 
    invece si sa che il prodotto scalare va calcolato sui punti della superficie $S$ e dunque boooh. 
    \begin{gather*}
        \int_{S}^{} (V_1 - V_2) (\vv{\nabla} \cdot V_1 \hat{n} - \vv{\nabla}V_2 \hat{n}    ) dS 
    \end{gather*}
    Si è trovato allora che 
    \begin{gather*}
        \int_{}^{}\left| \vv{\nabla} U  \right|^{2} \ d\tau = 0  
    \end{gather*}
    Dunque il gradiente è uguale a zero per tutti i punti del volume perché è tutta nulla:
    l'integrale di una funzione positiva ha integrale nullo se e solo se è nulla. Infatti la funzione 
    potenziale è definita positiva, dunque per tutti i punti del volume il gradiente
    di $U$ è zero.
    Per ogni punto si è trovato che
    \begin{gather*}
        \vv{\nabla}V_1 = \vv{\nabla}V_2  
    \end{gather*}
    Per un fisico il teorema è dimostrato poiché si è visto che il 
    campo è uguale. Tuttavia non è detto che il potenziale sia unico, infatti se
    $\vv{\nabla} U = 0$ allora è costante in tutti i punti di $\tau$. 
\end{proof}


\end{document}